
\section{Geometry in $\mathbb{R}^3$}

We will now start discussing geometry in $\mathbb{R}^3$.

Informally, $\mathbb{R}^3$ is three-dimensional, so inside it we can consider
two-dimensional flat objects, called planes, and one-dimensional flat objects,
called lines.

As we saw previously, it is easier to study these objects first in the case
that they pass through, or contain, the origin.

\subsection{Lines Through the Origin}

As in the case of $\mathbb{R}^2$, to describe a line through the origin, we
just need to give any nonzero vector along the line.

For example, let

\[
\vec v=
\begin{bmatrix}
3\\
4\\
5
\end{bmatrix}.
\]

The line through the origin in the direction of $\vec v$ is

\[
\operatorname{span}(\vec v)
=
\left\{
t
\begin{bmatrix}
3\\
4\\
5
\end{bmatrix}
:
t\in\mathbb{R}
\right\}.
\]

\begin{figure}[h]
\centering
\begin{tikzpicture}[
  scale=0.8,
  x={(0.85cm,-0.25cm)},
  y={(-0.35cm,-0.25cm)},
  z={(0cm,0.75cm)}
]
  % Origin
  \coordinate (O) at (0,0,0);

  % Coordinate axes
  \draw[->, gray] (0,0,0) -- (6,0,0) node[below right] {$x$};
  \draw[->, gray] (0,0,0) -- (0,6,0) node[below left] {$y$};
  \draw[->, gray] (0,0,0) -- (0,0,7) node[above] {$z$};

  % Exact line t(3,4,5)
  \draw[black, thick]
    (-2.4,-3.2,-4) -- (3.6,4.8,6);

  % Exact vector v=(3,4,5)
  \draw[-{Latex[length=3mm]}, very thick]
    (O) -- (3,4,5);

  \filldraw[black] (O) circle (1.7pt);

  \node[right=5pt] at (3,4,5)
    {$\vec v=\begin{bmatrix}3\\4\\5\end{bmatrix}$};

  \node[right=4pt] at (3.45,4.6,5.75) {$\ell$};
\end{tikzpicture}
\caption{The line through the origin spanned by
$\vec v=\begin{bmatrix}3\\4\\5\end{bmatrix}$.}
\end{figure}

\subsection{Planes Through the Origin}

Informally, a plane has two independent directions. This means that, to
describe a plane, we need to give two vectors in the plane that are not
multiples of each other.

For example, let

\[
\vec v_1=
\begin{bmatrix}
3\\
4\\
5
\end{bmatrix},
\qquad
\vec v_2=
\begin{bmatrix}
1\\
1\\
1
\end{bmatrix}.
\]

These vectors are not multiples of each other.

\begin{figure}[h]
\centering
\begin{tikzpicture}[
  scale=0.72,
  x={(0.85cm,-0.25cm)},
  y={(-0.35cm,-0.25cm)},
  z={(0cm,0.75cm)}
]
  \coordinate (O) at (0,0,0);

  % Coordinate axes
  \draw[->, gray] (0,0,0) -- (6,0,0) node[below right] {$x$};
  \draw[->, gray] (0,0,0) -- (0,6,0) node[below left] {$y$};
  \draw[->, gray] (0,0,0) -- (0,0,7) node[above] {$z$};

  % Exact plane x-2y+z=0
  % Each corner satisfies x-2y+z=0.
  \filldraw[
    fill=blue!15,
    fill opacity=0.65,
    draw=blue!55!black,
    thick
  ]
    (-4,-2,0) --
    (0,-2,-4) --
    (4,2,0) --
    (0,2,4) -- cycle;

  % Exact vectors
  \draw[-{Latex[length=3mm]}, very thick, black]
    (O) -- (3,4,5);
  \node[right=5pt] at (3,4,5)
    {$\vec v_1=\begin{bmatrix}3\\4\\5\end{bmatrix}$};

  \draw[-{Latex[length=3mm]}, very thick, black]
    (O) -- (1,1,1);
  \node[below right=4pt] at (1,1,1)
    {$\vec v_2=\begin{bmatrix}1\\1\\1\end{bmatrix}$};

  \filldraw[black] (O) circle (1.7pt);
  \node[right] at (3.5,1.75,0) {$P$};
\end{tikzpicture}
\caption{The vectors $\vec v_1$ and $\vec v_2$ determine a plane $P$ through
the origin.}
\end{figure}

The reason that two vectors are enough to describe a plane is that we can take
linear combinations of them to fill out the entire plane.

Let $\ell_1$ be the line through $\vec v_1$, and let $\ell_2$ be the line
through $\vec v_2$.

Consider a vector $\vec w$ in the plane. Complete a parallelogram whose sides
are parallel to $\ell_1$ and $\ell_2$. We can then write

\[
\vec w=\vec w_1+\vec w_2,
\]

where $\vec w_1$ lies on $\ell_1$ and $\vec w_2$ lies on $\ell_2$.

This implies that

\[
\vec w_1=a\vec v_1
\]

for some scalar $a$, and

\[
\vec w_2=b\vec v_2
\]

for some scalar $b$.

Therefore,

\[
\boxed{\vec w=a\vec v_1+b\vec v_2}.
\]

We say that $\vec w$ is a linear combination of $\vec v_1$ and $\vec v_2$.

The following figure gives an exact example. Here

\[
\vec w_1=\vec v_1,
\qquad
\vec w_2=-\vec v_2,
\]

so

\[
\vec w=\vec v_1-\vec v_2
=
\begin{bmatrix}
2\\
3\\
4
\end{bmatrix}.
\]

\begin{figure}[h]
\centering
\begin{tikzpicture}[
  scale=0.78,
  x={(0.85cm,-0.25cm)},
  y={(-0.35cm,-0.25cm)},
  z={(0cm,0.75cm)}
]
  \coordinate (O) at (0,0,0);
  \coordinate (Wone) at (3,4,5);
  \coordinate (Wtwo) at (-1,-1,-1);
  \coordinate (W) at (2,3,4);

  % Exact plane x-2y+z=0
  \filldraw[
    fill=blue!12,
    fill opacity=0.65,
    draw=blue!55!black,
    thick
  ]
    (-4,-2,0) --
    (0,-2,-4) --
    (4,2,0) --
    (0,2,4) -- cycle;

  % Lines in the directions v1 and v2
  \draw[blue!70!black, thick]
    (-2.4,-3.2,-4) -- (3.6,4.8,6);
  \node[right] at (3.3,4.4,5.5) {$\ell_1$};

  \draw[blue!70!black, thick]
    (-3,-3,-3) -- (3,3,3);
  \node[below right] at (2.5,2.5,2.5) {$\ell_2$};

  % Exact component vectors
  \draw[-{Latex[length=3mm]}, very thick]
    (O) -- (Wone);
  \node[right=5pt] at (Wone) {$\vec w_1=\vec v_1$};

  \draw[-{Latex[length=3mm]}, very thick]
    (O) -- (Wtwo);
  \node[left=5pt] at (Wtwo) {$\vec w_2=-\vec v_2$};

  % Exact sum vector
  \draw[-{Latex[length=3mm]}, very thick, green!55!black]
    (O) -- (W);
  \node[left=5pt, green!55!black] at (W)
    {$\vec w=\begin{bmatrix}2\\3\\4\end{bmatrix}$};

  % Exact parallelogram completion
  \draw[dashed, magenta!80!black] (Wone) -- (W);
  \draw[dashed, magenta!80!black] (Wtwo) -- (W);

  \filldraw[black] (O) circle (1.7pt);
  \node[right] at (3.5,1.75,0) {$P$};
\end{tikzpicture}
\caption{The exact parallelogram relation
$\vec w=\vec w_1+\vec w_2=\vec v_1-\vec v_2$.}
\end{figure}

\paragraph{Definition.}
The span of the two vectors $\vec v_1$ and $\vec v_2$ is the set of all linear
combinations

\[
\operatorname{span}(\vec v_1,\vec v_2)
=
\left\{
a\vec v_1+b\vec v_2
:
a,b\in\mathbb{R}
\right\}.
\]

Later, we will generalize this to the span of any number of vectors in
$\mathbb{R}^n$.

Returning to our specific example,

\[
\vec v_1=
\begin{bmatrix}
3\\
4\\
5
\end{bmatrix},
\qquad
\vec v_2=
\begin{bmatrix}
1\\
1\\
1
\end{bmatrix},
\]

the plane $P$ is described by

\[
P
=
\left\{
a\vec v_1+b\vec v_2
:
a,b\in\mathbb{R}
\right\}.
\]

Equivalently,

\[
P
=
\left\{
a
\begin{bmatrix}
3\\
4\\
5
\end{bmatrix}
+
b
\begin{bmatrix}
1\\
1\\
1
\end{bmatrix}
:
a,b\in\mathbb{R}
\right\}.
\]

Thus,

\[
P
=
\left\{
\begin{bmatrix}
3a+b\\
4a+b\\
5a+b
\end{bmatrix}
:
a,b\in\mathbb{R}
\right\}.
\]

\paragraph{Exercise.}
Use this last description to write out some other vectors in the plane.

\paragraph{Answer.}
We can take any values of $a$ and $b$.

For example, if $a=1$ and $b=-1$, then

\[
\begin{bmatrix}
3a+b\\
4a+b\\
5a+b
\end{bmatrix}
=
\begin{bmatrix}
2\\
3\\
4
\end{bmatrix}
=
\vec v_3.
\]

If $a=2$ and $b=3$, then

\[
\begin{bmatrix}
3a+b\\
4a+b\\
5a+b
\end{bmatrix}
=
\begin{bmatrix}
9\\
11\\
13
\end{bmatrix}
=
\vec v_4.
\]

Notice that taking $a=1$ and $b=0$ recovers $\vec v_1$, while taking $a=0$
and $b=1$ recovers $\vec v_2$.

Now there is nothing special about the original choice $(\vec v_1,\vec v_2)$.
As far as the plane $P$ is concerned, any choice of two independent vectors is
enough to describe the plane.

For example,

\[
P=\operatorname{span}(\vec v_3,\vec v_4).
\]

Why is this?

The key point is that $\vec v_1$ and $\vec v_2$ are themselves linear
combinations of $\vec v_3$ and $\vec v_4$.

Since

\[
\vec v_3=\vec v_1-\vec v_2
\]

and

\[
\vec v_4=2\vec v_1+3\vec v_2,
\]

we have

\[
\vec v_1
=
\frac{3\vec v_3+\vec v_4}{5}
=
\frac35\vec v_3+\frac15\vec v_4
\]

and

\[
\vec v_2
=
\frac{\vec v_4-2\vec v_3}{5}
=
-\frac25\vec v_3+\frac15\vec v_4.
\]

Therefore,

\[
\vec v_1,\vec v_2\in\operatorname{span}(\vec v_3,\vec v_4).
\]

The span of any number of vectors is closed under taking further linear
combinations. One might summarize this with the slogan:

\begin{center}
A linear combination of linear combinations is a linear combination.
\end{center}

Indeed, if

\[
\vec v_1=\alpha\vec v_3+\beta\vec v_4
\]

and

\[
\vec v_2=\gamma\vec v_3+\delta\vec v_4,
\]

then

\[
\begin{aligned}
a\vec v_1+b\vec v_2
&=
a(\alpha\vec v_3+\beta\vec v_4)
+
b(\gamma\vec v_3+\delta\vec v_4)\\
&=
(a\alpha+b\gamma)\vec v_3
+
(a\beta+b\delta)\vec v_4\\
&\in
\operatorname{span}(\vec v_3,\vec v_4).
\end{aligned}
\]

This means that any linear combination of $(\vec v_1,\vec v_2)$ is also a
linear combination of $(\vec v_3,\vec v_4)$.

Likewise, any linear combination of $(\vec v_3,\vec v_4)$ is a linear
combination of $(\vec v_1,\vec v_2)$.

Therefore,

\[
\operatorname{span}(\vec v_1,\vec v_2)
=
\operatorname{span}(\vec v_3,\vec v_4).
\]

\paragraph{Conclusion.}
A plane in $\mathbb{R}^3$ containing the origin may be described as the span
of two independent vectors. This description is not unique, since there are
many different choices of pairs of independent vectors.

\subsection{The Dual Perspective}

Now we will consider lines and planes in $\mathbb{R}^3$ from the dual
perspective. Instead of describing them as spans of vectors, we will try to
describe them using equations.

It is simplest to start with planes through the origin.

As an example, consider the plane $P$ discussed above:

\[
P=\operatorname{span}(\vec v_1,\vec v_2),
\]

where

\[
\vec v_1=
\begin{bmatrix}
3\\
4\\
5
\end{bmatrix},
\qquad
\vec v_2=
\begin{bmatrix}
1\\
1\\
1
\end{bmatrix}.
\]

We can guess an equation defining this plane by noticing that, for both
$\vec v_1$ and $\vec v_2$,

\[
\text{$x$-coordinate}+\text{$z$-coordinate}
=
2\cdot\text{$y$-coordinate}.
\]

Indeed,

\[
3+5=2(4)
\]

and

\[
1+1=2(1).
\]

This means that $\vec v_1$ and $\vec v_2$ satisfy the equation

\[
x+z=2y,
\]

or equivalently,

\[
x-2y+z=0.
\]

\paragraph{Exercise.}
Graph the equation

\[
x-2y+z=0.
\]

Is it the same as the plane $P$?

What is happening here?

Let

\[
\vec w=
\begin{bmatrix}
1\\
-2\\
1
\end{bmatrix}.
\]

\begin{figure}[h]
\centering
\begin{tikzpicture}[
  scale=0.8,
  x={(0.85cm,-0.25cm)},
  y={(-0.35cm,-0.25cm)},
  z={(0cm,0.75cm)}
]
  \coordinate (O) at (0,0,0);

  % Exact plane x-2y+z=0
  \filldraw[
    fill=blue!15,
    fill opacity=0.65,
    draw=blue!55!black,
    thick
  ]
    (-4,-2,0) --
    (0,-2,-4) --
    (4,2,0) --
    (0,2,4) -- cycle;

  % Exact spanning vectors
  \draw[-{Latex[length=3mm]}, very thick]
    (O) -- (3,4,5);
  \node[right=4pt] at (3,4,5) {$\vec v_1$};

  \draw[-{Latex[length=3mm]}, very thick]
    (O) -- (1,1,1);
  \node[below right=3pt] at (1,1,1) {$\vec v_2$};

  % Exact normal line span(1,-2,1)
  \draw[green!55!black, thick]
    (-2,4,-2) -- (2,-4,2);

  % Exact normal vector
  \draw[-{Latex[length=3mm]}, very thick, green!55!black]
    (O) -- (1,-2,1);
  \node[right=5pt, green!55!black] at (1,-2,1)
    {$\vec w=\begin{bmatrix}1\\-2\\1\end{bmatrix}$};

  \filldraw[black] (O) circle (1.7pt);
  \node[right] at (3.5,1.75,0) {$P$};
\end{tikzpicture}
\caption{The vector
$\vec w=\begin{bmatrix}1\\-2\\1\end{bmatrix}$ is perpendicular to the plane
$P$.}
\end{figure}

The vector

\[
\vec w=
\begin{bmatrix}
1\\
-2\\
1
\end{bmatrix}
\]

spans the unique line through the origin that is perpendicular to the plane
$P$.

Notice that

\[
\vec v_1\cdot\vec w
=
\begin{bmatrix}
3\\
4\\
5
\end{bmatrix}
\cdot
\begin{bmatrix}
1\\
-2\\
1
\end{bmatrix}
=
3-8+5
=
0,
\]

and

\[
\vec v_2\cdot\vec w
=
\begin{bmatrix}
1\\
1\\
1
\end{bmatrix}
\cdot
\begin{bmatrix}
1\\
-2\\
1
\end{bmatrix}
=
1-2+1
=
0.
\]

Consequently, for any linear combination of $\vec v_1$ and $\vec v_2$,

\[
\begin{aligned}
(a\vec v_1+b\vec v_2)\cdot\vec w
&=
a(\vec v_1\cdot\vec w)
+
b(\vec v_2\cdot\vec w)\\
&=
0+0\\
&=
0.
\end{aligned}
\]

\paragraph{Conclusion.}
The equation

\[
x-2y+z=0
\]

exactly describes the vectors perpendicular to

\[
\begin{bmatrix}
1\\
-2\\
1
\end{bmatrix}.
\]

This is another way of describing the plane $P$.

\paragraph{Remark.}
Instead of

\[
\begin{bmatrix}
1\\
-2\\
1
\end{bmatrix},
\]

we could take any nonzero scalar multiple of it. For example,

\[
\begin{bmatrix}
3\\
-6\\
3
\end{bmatrix}.
\]

Thus,

\[
3x-6y+3z=0
\]

is another equation describing $P$.

\paragraph{Question for in-class discussion.}
Which method of describing $P$ is easier:

\begin{itemize}
  \item as the span of two vectors, or
  \item as the solution set of a single linear equation?
\end{itemize}

\subsection{Describing Lines Using the Dual Perspective}

Now let us try to describe lines using the dual perspective.

This is a trickier problem, since one equation in $\mathbb{R}^3$ defines a
plane. To define a line, we need two equations.

For example, consider the line

\[
\ell
=
\operatorname{span}
\left(
\begin{bmatrix}
3\\
4\\
5
\end{bmatrix}
\right).
\]

One equation satisfied by every vector on this line is

\[
x-2y+z=0.
\]

Can we find another equation that is independent of the first equation?

There are many ways to do this. For example, each of the following equations
is satisfied by every vector on $\ell$:

\[
4x-3y=0,
\]

\[
5y-4z=0,
\]

\[
5x-3z=0,
\]

or

\[
-5x+5y-z=0.
\]

Let

\[
\vec w=
\begin{bmatrix}
1\\
-2\\
1
\end{bmatrix}
\]

and, just for the sake of an example, choose

\[
\vec w'=
\begin{bmatrix}
-5\\
5\\
-1
\end{bmatrix}.
\]

The plane perpendicular to $\vec w$ is

\[
P=\left\{\vec x\in\mathbb{R}^3:\vec x\cdot\vec w=0\right\},
\]

and the plane perpendicular to $\vec w'$ is

\[
P'=\left\{\vec x\in\mathbb{R}^3:\vec x\cdot\vec w'=0\right\}.
\]

Their intersection is the line $\ell$.

\begin{figure}[h]
\centering
\begin{tikzpicture}[
  scale=0.72,
  x={(0.85cm,-0.25cm)},
  y={(-0.35cm,-0.25cm)},
  z={(0cm,0.75cm)}
]
  \coordinate (O) at (0,0,0);

  % Plane P: span((3,4,5),(2,1,0))
  % Its normal vector is (1,-2,1).
  \filldraw[
    fill=blue!25,
    fill opacity=0.42,
    draw=blue!60!black,
    thick
  ]
    (-6.4,-5.2,-4) --
    (1.6,-1.2,-4) --
    (6.4,5.2,4) --
    (-1.6,1.2,4) -- cycle;

  % Plane P': span((3,4,5),(1,1,0))
  % Its normal vector is (-5,5,-1).
  \filldraw[
    fill=orange!35,
    fill opacity=0.40,
    draw=orange!70!black,
    thick
  ]
    (-4.4,-5.2,-4) --
    (-0.4,-1.2,-4) --
    (4.4,5.2,4) --
    (0.4,1.2,4) -- cycle;

  % Exact intersection line span(3,4,5)
  \draw[black, very thick]
    (-3,-4,-5) -- (3.6,4.8,6);

  % Exact direction vector
  \draw[-{Latex[length=3mm]}, very thick, black]
    (O) -- (3,4,5);
  \node[right=5pt] at (3,4,5)
    {$\begin{bmatrix}3\\4\\5\end{bmatrix}$};

  % Exact normal vectors
  \draw[-{Latex[length=3mm]}, thick, blue!70!black]
    (O) -- (1,-2,1);
  \node[right=4pt, blue!70!black] at (1,-2,1)
    {$\vec w=\begin{bmatrix}1\\-2\\1\end{bmatrix}$};

  \draw[-{Latex[length=3mm]}, thick, orange!80!black]
    (O) -- (-5,5,-1);
  \node[left=5pt, orange!80!black] at (-5,5,-1)
    {$\vec w'=\begin{bmatrix}-5\\5\\-1\end{bmatrix}$};

  \filldraw[black] (O) circle (1.7pt);

  \node[blue!70!black] at (4.2,3.5,1.5) {$P$};
  \node[orange!80!black] at (-2.6,-3.2,-1.7) {$P'$};
  \node[right] at (3.4,4.55,5.7) {$\ell$};
\end{tikzpicture}
\caption{The line
$\ell=\operatorname{span}\left(\begin{bmatrix}3\\4\\5\end{bmatrix}\right)$
is the intersection of two planes through the origin.}
\end{figure}

Thus, the equations

\[
\begin{cases}
x-2y+z=0,\\
-5x+5y-z=0
\end{cases}
\]

describe the line $\ell$.

\paragraph{Remark.}
Each equation defines a plane. What we are doing is writing $\ell$ as the
intersection of two planes:

\[
\ell
=
\left\{
\vec x\in\mathbb{R}^3:
\vec x\cdot\vec w=0
\text{ and }
\vec x\cdot\vec w'=0
\right\}.
\]

There are many different ways to do this. Each such way corresponds to choosing
a pair of independent vectors $(\vec w_1,\vec w_2)$ in the plane perpendicular
to $\ell$.

\subsection{Affine Lines and Planes}

Affine lines and planes are obtained by translating a line or plane through the
origin by a fixed vector.

An affine line is given by

\[
\vec v_0+\operatorname{span}(\vec v_1)
\]

for some vectors $\vec v_0$ and $\vec v_1$, with $\vec v_1\ne\vec 0$.

An affine plane is given by

\[
\vec v_0+\operatorname{span}(\vec v_1,\vec v_2)
\]

for two independent vectors $\vec v_1$ and $\vec v_2$.

Dually, in terms of equations, an affine plane is given by a single
nonhomogeneous linear equation

\[
ax+by+cz=d,
\]

where $(a,b,c)\ne(0,0,0)$.

An affine line in $\mathbb{R}^3$ is given by two independent linear equations,

\[
\begin{cases}
a_1x+b_1y+c_1z=d_1,\\
a_2x+b_2y+c_2z=d_2,
\end{cases}
\]

whose corresponding normal vectors

\[
\begin{bmatrix}
a_1\\
b_1\\
c_1
\end{bmatrix},
\qquad
\begin{bmatrix}
a_2\\
b_2\\
c_2
\end{bmatrix}
\]

are independent, meaning that they are not scalar multiples of each other.